\chapter{The spectral reading}

This chapter reads the dependency matrix with every cell and not its largest. That finds a law the maximum cannot see, and then confirms the null harder than the maximum can.

\section{Why a maximum is the wrong reading}

Every strict avalanche reading in the literature takes the largest of \(131{,}072\) cells and discards the other \(131{,}071\). Discarding them pins the peak at \(\sqrt{2 \ln N}\) whatever the function does. A structure spread thinly across many cells is invisible to a maximum and plain to a sum.

SHA-256 is built from rotations, and a rotation makes bit \(j\) depend on bit \(j - r\). If any of that survives into the dependency matrix it appears as a function of \((\text{input bit} - \text{output bit}) \bmod 32\), spread over every word pair and concentrated nowhere. Folding the cells onto that residue puts \(4096\) of them behind each number instead of one.

The fold is signed and not absolute, deliberately. Taking magnitudes first would fold noise into a positive bias and manufacture a signal in every class.

\section{What all thirty-two classes share}
\label{sec:commonmode}

Two things put a large offset into \emph{every} class at once. The avalanche is incomplete at these depths, and most cells sit far from one half. And the message schedule imposes a light cone: round \(t\) consumes \(W_t\). Message word \(w\) cannot affect the state before round \(w+1\), and until then flipping any of its bits changes nothing at all, a cell pinned at \(-\sqrt{n}\) exactly. Neither of those depends on the output bit. Both spread evenly over all thirty-two residues and concentrate in none of them.

The fold therefore has two components, and the class mean at each depth separates them: a common mode every class shares, and a class pattern measured from that mean.

\subsection{The common mode decays at one thirty-second per round}

The loudest class, read without subtracting the class mean, tracks the common mode with a fixed offset. Its amplitude in standard errors over rounds twelve to twenty-three:

\begin{center}
\begin{tabular}{rrrl}
\toprule
rounds & loudest residue & sigmas & reading \\
\midrule
12 & 26 & 15773.15 & common mode \\
13 & 26 & 13724.82 & common mode \\
14 & 26 & 11676.61 & common mode \\
15 & 26 & 9627.78 & common mode \\
16 & 26 & 7580.54 & common mode \\
17 & 26 & 5546.15 & common mode \\
18 & 26 & 3551.43 & common mode \\
19 & 26 & 1916.81 & common mode \\
20 & 26 & 789.13 & common mode \\
21 & 26 & 160.61 & common mode \\
22 & 1 & 13.13 & collapsing \\
23 & 12 & 2.09 & flat \\
32 & 8 & 2.07 & flat \\
64 & 31 & 2.04 & flat \\
\bottomrule
\end{tabular}
\end{center}

The first four differences are \(2048.33\), \(2048.21\), \(2048.83\) and \(2047.24\). One sigma of the class mean at that sample is \(1.526 \times 10^{-5}\). \(2048\) sigmas is exactly \(1/32\).

\textbf{The common mode loses one thirty-second of its amplitude per round}, one bit position of thirty-two going diffuse, holding to four decimal places for five rounds, then accelerating into collapse at twenty-two. It is a quantized linear decay and not a smooth exponential, as a function made of discrete operations should produce and as a continuous model would not predict.

\textbf{Residue twenty-six is not a ridge.} Its deviation from the class mean is \(-1.36\) sigmas at rounds eight, twelve and sixteen alike, against a null peak of \(2.63\) for the largest of thirty-two, and it is a deficit:

\begin{center}
\begin{tabular}{rrrrrr}
\toprule
round & raw at 26 & class mean & excess & scatter & sigmas \\
\midrule
8 & \(-23964.8\) & \(-22049.6\) & \(-1915.2\) & 1410.9 & \(-1.36\) \\
12 & \(-15773.8\) & \(-13857.8\) & \(-1916.0\) & 1411.2 & \(-1.36\) \\
16 & \(-7580.8\) & \(-5665.9\) & \(-1914.9\) & 1411.0 & \(-1.36\) \\
\bottomrule
\end{tabular}
\end{center}

A reading that takes the loudest class and never subtracts the class mean reports a fixed ridge at residue twenty-six decaying at \(1/32\) per round. The decay is real and belongs to the common mode.

\subsection{The law at four times the sample}

The law comes out of a sample where \(2048\) sigmas happens to equal one thirty-second, and \(2048\) is \(2^{11}\) while the trials are \(2^{18}\). Every clean power of two in that arithmetic comes from the chosen sample size, and a law exact at only one sample size is the shape of the measurement and not a law.

\begin{center}
\begin{tabular}{lrrr}
\toprule
 & \(2^{18}\) trials & \(2^{20}\) trials & ratio \\
\midrule
round 12 & 15773.15 & 31546.74 & 2.0001 \\
round 13 & 13724.82 & 27450.50 & \\
round 14 & 11676.61 & 23354.13 & \\
\midrule
differences & 2048.33 & 4096.24 & \\
 & 2048.21 & 4096.37 & \\
\bottomrule
\end{tabular}
\end{center}

Sigma doubles for four times the trials, which is \(\sqrt{k}\) exactly, while the absolute deviation holds at \(4096 \times 2^{-17} = 1/32\). That is the signature of signal, and the precise opposite of the refuted avalanche dip, which holds its sigma and shrinks its deviation.

\subsection{The class pattern does not move}

The fold with the common mode removed is not empty, and what it holds is the round function's own bit transport read straight off the matrix. The ranking is identical at every depth from six to nineteen.

\begin{center}
\begin{tabular}{rlll}
\toprule
rank & class & sigmas at round 12 & what it is \\
\midrule
1 & 0 & \(+3.95\) & the diagonal: bit \(j\) drives bit \(j\) \\
2 & 25 & \(+1.84\) & \(\Sigma_1\), \(\mathrm{ROTR}25\) \\
3 & 6 & \(+1.60\) & \(\Sigma_1\), \(\mathrm{ROTR}6\) \\
4 & 31 & \(+1.38\) & \(-1 \bmod 32\): the carry, bit \(k\) reaching bit \(k+1\) \\
5 & 11 & \(+1.31\) & \(\Sigma_1\), \(\mathrm{ROTR}11\) \\
6 & 22 & \(+0.64\) & \(\Sigma_0\), \(\mathrm{ROTR}22\) \\
\bottomrule
\end{tabular}
\end{center}

Residue zero holds \(+3.95\) sigmas and is constant to four significant figures across eleven rounds, from six to sixteen, while the common mode falls by a factor of nearly three. The scatter of the other classes is constant at \(1174\) over the same span. The field is two components: a common mode that decays and a class pattern that does not move at all. Nothing rotates, and there is no moving structure for an inverse rotation to unwind.

Classes zero and thirty-one are predictable in advance: any account of this function expects the diagonal and the carry to dominate. \textbf{The finding is that \(\Sigma_1\)'s three amounts appear and \(\Sigma_0\)'s barely do.} All three of \(\{6, 11, 25\}\) land in the top six at every depth, about one chance in \(225\) at a single depth, or one in \(112\) allowing that \(\Sigma_0\) would have counted equally.

The mechanism is the asymmetry of the two halves. \(\Sigma_1\) acts on \(e\) and feeds \(T_1\), which is added into both the \(a\) chain and the \(e\) chain. \(\Sigma_0\) acts on \(a\) and feeds \(T_2\), which reaches the \(a\) chain only. Twice the reach, and it is visible in the spectrum.

This also explains why majority measures as contributing nothing on four separate tests. \textbf{Majority and \(\Sigma_0\) are both in \(T_2\)}, and \(T_2\) is the weak side. The mixing of SHA-256 is in \(T_1\).

\section{The strongest null in this work}

Past round twenty-three the fold reads \(2.09\), \(2.07\) and \(2.04\) against a null peak of \(2.63\), at three depths, with \(4096\) cells behind each number and not one. That is roughly sixty-four times the sensitivity the maximum has, finding nothing.

The harmonic sweep agrees. The round-twenty-three profile wanders within two sigmas with no visible pattern, and its discrete transform peaks at \(f_{13}\) with magnitude \(1.80\) against a null peak of \(2.35\): no repeating trend, no harmonics, no side lobes.

\section{What the spectral reading settles}

\textbf{The common mode decays by a law}, one thirty-second per round from round twelve to twenty-one, verified at two sample sizes.

\textbf{The class pattern is fixed}: the diagonal, \(\Sigma_1\)'s three amounts and the carry, ranked identically from round six to nineteen.

\textbf{Both are gone by round twenty-three}, at sixty-four times the sensitivity of a maximum.
